\newpage
\section{Dienste und Modelle}
  \subsection{Grundmodell der Kommunikation}
    \begin{itemize}
      \item Ein abstraktes Medium überbrückt die räumliche Distanz zwischen Nutzern (Sender, Empfänger).
        Es abstrahiert von der konkreten Umsetzung.
      \item Stellt einen (Kommunikations-)Dienst zur Verfügung
    \end{itemize}
  \subsection{Dienste}
    Dienstesicht
    \begin{itemize}
      \item Modellierung eines Rechnernetztes bzw. des abstrakten Mediums als Black-Box
        \begin{itemize}
          \item Abstraktes Medium ist Diensterbringer, erbringt einen Dienst für dessen Nutzer
          \item Nutzer sind Dienstnehmer: z.B. Prozess auf einem Rechner, Mensch
        \end{itemize}
    \end{itemize}
    Dienst
    \begin{itemize}
      \item Ein Dienst bündelt zusammengehörende Funktionen und stellt diese einem Nutzer (Dienstnehmer)
        zur Verfügung
        \begin{itemize}
          \item Einzelne Teile eines Dienstes können unabhängig voneinander in Anspruch genommen werden
        \end{itemize}
      \item Dienstnehmer nimmt ein Dienst in Anspruch
      \item Diensterbringer stellt einen Dienst zur Verfügung
      \item Dienstzugangspunkt (Service Acess Point, SAP) stellt die Schnittstelle zu einem Dienst dar
      \item Dienstprimitive beschreiben die Interaktion zwischen Dienstnehmer und Diensterbringer
        am Dienstzugangspunkt
      \item Die Beschreibung eine Dienstes reflektiert das Verhalten an den Dienstzugangspunkten zum
        Diensterbringer
    \end{itemize}
    Dienstzugangspunkt
    \begin{itemize}
      \item Typen von Dienstprimitiven (Schnittstellenereignissen)
        \begin{itemize}
          \item Request (req)
            \begin{itemize}
              \item Beauftragen (Dienstnehmer $\ra$ Diensterbringer)
            \end{itemize}
          \item Indication (ind)
            \begin{itemize}
              \item Benachrichtigung des Partners (Diensterbringer $\ra$ Dienstnehmer)
            \end{itemize}
          \item Response (rsp)
            \begin{itemize}
              \item Beantwortung durch Partner (Dienstnehmer $\ra$ Diensterbringer)
            \end{itemize}
          \item Confirmation (cnf)
            \begin{itemize}
              \item Benachrichtigung über Abschluss (Diensterbringer $\ra$ Dienstnehmer)
            \end{itemize}
        \end{itemize}
    \end{itemize}
    Grundformen von Diensten
    \begin{itemize}
      \item Unzuverlässiger Dienst
        \begin{itemize}
          \item Unbestätigter Dienst (Best Effort)
          \item Bestätigter Dienst
        \end{itemize}
      \item Zuverlässiger Dienst
    \end{itemize}
  \subsection{Schichtenmodell / Referenzmodell}
    Vereinfachung des Grundmodells
    \begin{itemize}
      \item Prinzip
        \begin{itemize}
          \item Einführung einer Abstraktion auf der Basis von Schichten
            \begin{itemize}
              \item Bereitstellen von Diensten an der (Dienst-)Schnittstelle zum Dienstnehmer
              \item Nutzung von Diensten an der Schnittstelle nach unten
            \end{itemize}
        \end{itemize}
    \end{itemize}
    Schicht
    \begin{itemize}
      \item Eine Schicht bündelt zusammengehörende Funktionalitäten, die sich von denen anderer
        Schichten unterscheiden
      \item Eine Schicht bietet an ihrer Dienstschnittstelle wohldefinierte Dienste an     
        \begin{itemize}
          \item Die technische Umsetzung der Schicht ist an dieser Dienstschnittstelle nicht sichtbar
            (transparent, opak) $\ra$ Black-Box
          \item Zur Erbringung dieser Dienste nutzt eine Schicht die Dienste der darunter liegenden
            Schicht
        \end{itemize}
      \item Schichten sind hierarchisch angeordnet
        \begin{itemize}
          \item Schichten können nicht entfallen
        \end{itemize}
      \item Eine Schicht kann weiter in Teilschichten geglieder werden
        \begin{itemize}
          \item Teilschichten können entfallen
        \end{itemize}
    \end{itemize}
    Protokolle und Dienste
    \begin{itemize}
      \item Protokolle definieren Regeln und Formate für die Kommunikation (Interaktion) zwischen zwei
        oder mehr Instanzen sowie die beim Senden bzw. Empfangen von Daten bzw. Ereignissen jeweils
        durchzuführenden Aktionen
      \item Protokolle sind innerhalb einer Schicht lokalisiert
        \begin{itemize}
          \item Werden durch Protokollinstanzen ausgeführt
          \item Protokollinstanz ist i.d.R. einem Gerät zugeordnet
        \end{itemize}
      \item Dienst einer Schicht wird durch Zusammenwirken der Protokollinstanzen dieser Schicht
        erbracht
        \begin{itemize}
          \item Schicht zieht sich über mehrer Geräte hinweg. Sie ist nicht auf ein einzelnes Gerät
            limitiert
        \end{itemize}
    \end{itemize}
    Horizontale und vertikale Kommunikation
    \begin{itemize}
      \item Horizontale Kommunikation
        \begin{itemize}
          \item Zwischen Sender und Empfänger (innerhalb einer Schicht)
          \item Protokollinstanzen einer Schicht tauschen Daten untereinander aus um den geforderten
            Dienst zu erbringen
        \end{itemize}
      \item Vertikale Kommunikation
        \begin{itemize}
          \item Zwischen Schichten (innerhalb eines Gerätes)
          \item Protokollinstanz der Schicht N greift auf Dienste der Protokollinstanz in Schicht
            N-1 zu bzw. gibt Daten an Protokollinstanz der Schicht N+1 weiter
        \end{itemize}
    \end{itemize}
    Physikalische Schicht (Physical Layer)
    \begin{itemize}
      \item Zwei direkt über ein physikalisches Übertragungsmedium verbundene Geräte können Bits
        austauschen
      \item Aufgaben
        \begin{itemize}
          \item Übertragung von (unstrukturierten) Bitfolgen über ein längenbeschränktes physikalisches
            Medium
        \end{itemize}
      \item Eigenschaften
        \begin{itemize}
          \item Bietet Schicht 2 einen unzuverlässigen Dienst an
          \item Verwendung von Leitungscodes, Modulation etc.
          \item Auf dem Medium werden Bits nicht gepuffert
          \item Bei einer Übertragung können Störungen auftreten
        \end{itemize}
    \end{itemize}
    Sicherungsschicht (Data Link Layer)
    \begin{itemize}
      \item Übertragung von Paketen zwischen physikalisch benachbarten Geräten
        \begin{itemize}
          \item Strukturierte Übertragung
          \item Pakete werden hier als Rahmen bezeichnet
        \end{itemize}
      \item Aufgabe
        \begin{itemize}
          \item Kommunikation zwischen physikalisch benachbarten Systemen
        \end{itemize}
      \item Eigenschaften
        \begin{itemize}
          \item Bereitstellung zuverlässiger und unzuverlässiger Dienste
            \begin{itemize}
              \item Zur Bereitstellung zuverlässiger Dienste Erkennung und Behebung von Fehlern der
                physikalischen Schicht
            \end{itemize}
          \item Adressierung der Geräte
          \item Pufferung der Daten sowohl beim Sender als auch beim Empfänger
          \item Strukturierung der Übertragung in Rahmen
          \item Abstrahierung vonm physikalischen Medium
        \end{itemize}
    \end{itemize}
    Vermittlungsschicht (Network Layer)
    \begin{itemize}
      \item Datenaustausch zwischen Endsystemen
        \begin{itemize}
          \item Sind i.d.R. nicht physikalisch benachbart
          \item Sind nicht mit gleichem physikalischen Medium verbunden
            \begin{itemize}
              \item Benötigt Zwischensysteme
            \end{itemize}
          \item Pakete werden hier typischerweise als Datagramme bezeichnet
        \end{itemize}
      \item Aufgabe
        \begin{itemize}
          \item Verknüpft Teilstrecken zwischen Geräten zu Ende-zu-Ende-Strecken
        \end{itemize}
      \item Eigenschaften
        \begin{itemize}
          \item Wegewahl
            \begin{itemize}
              \item Finden geeigneter Wege zwischen kommuniziernden Endsystemen
            \end{itemize}
          \item Adressierung von Endsystemen
          \item Bereitstellung zuverlässiger und unzuverlässiger Dienste
          \item Datagramme können in Zwischensystemen gepuffert werden
          \item Datagramme der Vermittlungsschicht werden in einem oder mehreren Rahmen der
            Sicherungsschicht übertragen
        \end{itemize}
    \end{itemize}
    Tranpsortschicht (Transport Layer)
    \begin{itemize}
      \item Datenaustausch zwischen Anwendungsprozessen auf Endsystemen
        \begin{itemize}
          \item Pakete werden hier typischerweise als Segmente oder Datagramme bezeichnet
        \end{itemize}
      \item Aufgabe
        \begin{itemize}
          \item Übertragung von Daten zwischen Anwendungsprozessen
        \end{itemize}
      \item Eigenschaften
        \begin{itemize}
          \item Bereitstellung zuverlässiger und unzuverlässiger Dienste
            \begin{itemize}
              \item Fehlererkennung und -behebung für zuverlässige Dienste
            \end{itemize}
          \item Pufferung der Daten in den Endsystemen
          \item Adressierung von Anwendungsprozessen
          \item Segmente/Datagramme der Transportschicht werden in ein oder mehrere Datagrammen der
            Vermittlungsschicht übertragen
        \end{itemize}
    \end{itemize}
    Nutzer-zu-Nutzer vs Ende-zu-Ende
    \begin{itemize}
      \item Transportschicht
        \begin{itemize}
          \item Logische Kommunikation zwischen (Anwendungs-)Prozessen
          \item Nutzer-zu-Nutzer
          \item Benutzt hierzu die Dienste der Vermittlungsschicht
        \end{itemize}
      \item Vermittlungsschicht
        \begin{itemize}
          \item Logische Kommunikation zwischen Endsystemen
          \item Ende-zu-Ende
          \item Benutzt hierzu die Dienste der Sicherungsschicht
        \end{itemize}
    \end{itemize}
    Anwendungsschicht
    \begin{itemize}
      \item Logik der Anwendung und hierfür erforderlichen Mechanismen und Datenstrukturen
        \begin{itemize}
          \item Instanzen der Anwendungsschicht tauschen Nachrichten aus
        \end{itemize}
    \end{itemize}
    Zusammenspiel der Sichten
    \begin{itemize}
      \item Daten senden
        \begin{itemize}
          \item Erstellen und hinzufügen von Metainformationen (Kopf, Anhang)
        \end{itemize}
        \includegraphics[width=.9\textwidth]{img/schichten_senden.png}
      \item Daten empfangen
        \begin{itemize}
          \item Interpretieren und Löschen von Metainformationen (Kopf, Anhang)
        \end{itemize}
        \includegraphics[width=.9\textwidth]{img/schichten_empfangen.png}
    \end{itemize}
  \subsection{OSI-Referenzmodell}
    \begin{itemize}
      \item Von ISO standardisiert
        \begin{itemize}
          \item Umfasst sieben Schichten
          \item Häufig als "das" Referenzmodell gesehen
        \end{itemize}
        \includegraphics[width=.6\textwidth]{img/osi.png}
    \end{itemize}
    Transportorientierte Schichten
    \begin{itemize}
      \item Übertragung von Daten zwischen Prozessen
        \begin{itemize}
          \item Nur den Bedürfnissen des Datenaustauschs und des bereitzustellenden Dienstes unterstellt
          \item Funktionsweise unabhängig von Inhalt der Daten
        \end{itemize}
      \item Begriff Transportsystem: die transportorientierten Schichten
        \begin{itemize}
          \item Schichten 1 bis 4 identisch mit 5-Schichten-Modell
        \end{itemize}
    \end{itemize}
    Anwendungsorientierte Schichten
    \begin{itemize}
      \item Datenaustausch zwischen Anwendungsprozessen
      \item Anwendungsaspekte relevant
        \begin{itemize}
          \item Informationsdarstellung und Informationsaustausch
          \item Steuerung des Ablaufs nach anwendungsorientierten Gesichtspunkten
        \end{itemize}
        \item Nur in Systemene vorhanden, in denen sich Anwendungsinstanzen befinden
        \begin{itemize}
          \item Typischerweise Endsysteme
        \end{itemize}
    \end{itemize}
